\section{Introduction}
\label{sec:intro}

Packages for millimetre-wave front ends, antenna-in-package modules and wide-bandgap power devices now have to route conductors in three dimensions through a low-loss dielectric, and to do it close to the die \cite{lau2022advanced,watanabe2021review,zhang2019antenna,chen2017sic}. Co-fired ceramics answer part of this demand. Low-temperature co-fired ceramic (LTCC) and glass-ceramic substrates combine low dielectric loss, hermeticity and a thermal expansion close to that of silicon \cite{sebastian2008low,tummala1991ceramic}. They are, however, built from stacked tapes: conductors are printed in planes, interconnections run vertically through punched vias, and every change of geometry requires new tooling \cite{imanaka2005multilayered,sebastian2008low}. The Heterogeneous Integration Roadmap lists additive manufacturing of ceramic packages, including lithography-based ceramic manufacturing, as one way past these limits \cite{ieeeeps2024hir}.

The conductor of choice in fired packages has long been silver, because it can be co-fired in air. Silver migrates under bias and humidity and diffuses into the glass when fired in oxidising conditions \cite{yang2007failure,jean2004interfacial,ma2016suppression}. Copper conducts almost as well, costs less and is far less prone to migration, but it must be fired in an atmosphere that keeps it metallic. In such an atmosphere the organic binder no longer burns off cleanly: part of it is left as char, and char trapped in a glass that has already sealed its pores degrades dielectric loss and can bloat the part \cite{lewis1997binder}. IBM solved this for its glass-ceramic/copper multilayer substrates with a steam-based co-sintering process that gave a dense glass-ceramic with high-conductivity copper \cite{master1991cosintering,tummala1991ceramic}. That solution was developed for laminated tapes and its operating window was found empirically.

Vat photopolymerisation of ceramic slurries, in particular lithography-based ceramic manufacturing (LCM), produces dense ceramics with features of a few tens of micrometres \cite{schwentenwein2015additive,stampfl2023lithography,halloran2016ceramic}. It has already been used for monolithic microwave filters and antenna substrates \cite{qian2022printed,fontana2023novel,mazingue2021printed}, and photocurable LTCC suspensions have been formulated for it \cite{fernandes2021study}. In these parts the metallisation is still added afterwards, by plating or by printing on the fired ceramic \cite{qian2022printed,wang2026fabrication}, or the conductors are jetted into ceramic layers by a separate route \cite{hirao2019novel,jager2024inkjet,liang2023additive,duan2023costeffective}. Multi-material vat photopolymerisation removes that restriction. Printers such as the CeraFab Multi 2M30 switch between slurries within a layer \cite{geier2020multimaterial,xu2026vat,subedi2024multimaterial}. Co-sintered ceramic--ceramic parts have been demonstrated \cite{schlacher2024towards}, as has a zirconia part with titanium-oxide conductors co-fired in a reducing atmosphere \cite{schwarzerfischer2021combining}. Copper slurries for vat photopolymerisation have been developed separately \cite{roumanie2021influence,scheibler2024lithography,resch2023lithography,canocano2024recent}. In principle, then, a glass-ceramic body and its copper wiring can be printed as one green part. We found no study of whether such a part can be fired.

Four requirements make that firing difficult, and they are coupled. First, the char from both slurries must be removed before the glass seals its pores, in an atmosphere that does not oxidise the copper. Second, the two materials must shrink at compatible times and by compatible amounts; otherwise the part cambers, delaminates or cracks around its conductors \cite{cai1997constrainedI,green2008constrained,jean1997cofiring}. Third, the glass must densify before it crystallises, because crystals stop viscous flow. Fourth, a thermal-expansion mismatch of about 14~ppm\,K$^{-1}$ between copper and a cordierite-type glass-ceramic must be accommodated on cooling. In tape-based LTCC the second requirement is met by trial, most often by adding ceramic powder to the metal paste \cite{chang1998effects} or by adjusting the heating rate \cite{brandt2013improved}. Continuum models of co-sintering reproduce camber and mismatch stresses in ceramic and metal--ceramic bilayers \cite{olevsky2013sintering,molla2014finite,kanters2001cosintering,reiterer2006arrhenius,chretien2022distortion,antou2024cosintering}, but none includes the chemistry of binder removal, and none treats copper with a crystallising glass. The first and third requirements are therefore usually handled separately from the second, although in practice they compete for the same temperature range.

This paper sets out a model-based design of that firing for a copper conductor printed inside a cordierite-type glass-ceramic body. We extend a verified slab model of binder removal and sintering, originally developed for printed copper, to a crystallising glass-ceramic, and couple the two materials through the viscous analogue of laminate theory and a three-dimensional finite-element model built on the printed voxel grid. The model is verified against closed-form solutions and parameterised from published data, with every input tagged by its provenance. With it we (i) map the atmosphere and temperature window in which carbon can be removed from the glass before its pores close while copper stays metallic; (ii) show that fine copper and a fragile glass densify in different temperature ranges and with different curve shapes, and quantify the consequences for camber and conductor damage; (iii) design the copper paste and the firing programme together to bring the two into step, both at nominal inputs and robustly against their uncertainty, and follow the designs through a printed test vehicle in three dimensions; and (iv) identify the few material parameters that decide the outcome, and should therefore be measured first.
